\chapter{Минимальная конечная алгебра наблюдаемого baseline}

Проверим, какой минимальный Artin--Wedderburn menu способен породить нужные
неабелевы факторы и оставить один абелев генератор после unimodularity.
Рассматриваются блоки
\begin{equation}
\mathbb C,\quad \mathbb H,\quad M_2(\mathbb C),\quad
M_3(\mathbb C),\quad M_3(\mathbb R).
\end{equation}

\begin{proposition}[Минимум quaternionic branch]
В объявленном меню минимальная по вещественной размерности алгебра, которая
содержит quaternionic weak block, complex color block и два центральных
\(u(1)\)-направления до одной unimodularity constraint, равна
\begin{equation}
\boxed{\mathcal A_F^{\rm SM}=\mathbb C\oplus\mathbb H\oplus M_3(\mathbb C).}
\end{equation}
Её вещественная размерность равна \(2+4+18=24\).
\end{proposition}

До central constraint unitary Lie algebra имеет вид
\begin{equation}
u(1)_{\mathbb C}\oplus su(2)_{\mathbb H}\oplus
su(3)_{M_3}\oplus u(1)_{M_3}.
\end{equation}
Одна невырожденная связь между центральными генераторами оставляет
\begin{equation}
\boxed{su(3)_c\oplus su(2)_L\oplus u(1)_Y.}
\end{equation}

Без quaternionic condition существует baseline
\(M_2(\mathbb C)\oplus M_3(\mathbb C)\) размерности 26. Поэтому
минимальность является условной внутри quaternionic branch, а не абсолютной
классификационной теоремой.

\section{Bimodule target-ledger}

Минимальная target-реализация одного поколения использует рёбра
\begin{center}
\begin{tabular}{cccc}
\toprule
Multiplet & Левый блок & Правый блок & Размерность \\
\midrule
\(Q_L\) & \(\mathbb H\) & \(M_3(\mathbb C)\) & 6 \\
\(L_L\) & \(\mathbb H\) & \(\mathbb C\) & 2 \\
\(u_R\) & \(\mathbb C\) & \(M_3(\mathbb C)\) & 3 \\
\(d_R\) & \(\overline{\mathbb C}\) & \(M_3(\mathbb C)\) & 3 \\
\(\nu_R\) & \(\mathbb C\) & \(\mathbb C\) & 1 \\
\(e_R\) & \(\overline{\mathbb C}\) & \(\mathbb C\) & 1 \\
\bottomrule
\end{tabular}
\end{center}
Итого: 16 complex particle states до antiparticle doubling.

\begin{theorem}[Observed gauge-algebra baseline]
Внутри предзарегистрированного quaternionic menu алгебра
\(\mathbb C\oplus\mathbb H\oplus M_3(\mathbb C)\) является минимальным
baseline, чья unitary Lie algebra после одной central constraint имеет
локальную структуру Стандартной модели.
\end{theorem}

\begin{nogocriterion}[Baseline ещё не является происхождением]
Поскольку bimodule ledger пока выбран по известным частицам, доказаны
совместимость и ограниченная минимальность, но не уникальное происхождение
спектра. Открыты полный bimodule-перебор, discrete quotient, три поколения и
Yukawa matrices.
\end{nogocriterion}

Машинный результат сохранён в
\path{s2t_v4_finite_algebra_menu_gate_results.json}.